% !PRINTED 288
\weberpage{288}
%<<CONT
\weq{16}{\begin{pmatrix} \alpha, & \beta \\ \gamma, & \delta \end{pmatrix}, \quad \alpha\delta - \beta\gamma \equiv 1 \ (\mathrm{mod}\ p).}
In einer dieser Substitutionen sind die Zahlen $\alpha$, $\beta$, $\gamma$, $\delta$ nach dem Modul $p$ bis auf das gemeinsame Vorzeichen bestimmt. Der Grad der Gruppe $\mathfrak{L}_0$ ist
\weq{17}{\frac{1}{2}p(p^2-1).}

Auf die Form (16) kommt man aber direkt, wenn man als die Gruppe der Teilungsgleichung nicht die Gruppe $\mathfrak{A}$, sondern die Gruppe $\mathfrak{B}$ des § 63 betrachtet, d.\,h. wenn man \gesperrt{$p$te Einheitswurzeln dem Rationalitätsbereich adjungiert}, woraus der Satz fließt:

\gesperrt{Die Gruppe $\mathfrak{L}_0$ ist die Gruppe der Transformationsgleichung, wenn $p$te Einheitswurzeln dem Rationalitätsbereich adjungiert sind.}

Zur Reduktion der Gruppe $\mathfrak{L}$ auf die Gruppe $\mathfrak{L}_0$ genügt aber schon die Adjunktion einer zweiwertigen Funktion und daher ist mit der Adjunktion der $p$ten Einheitswurzeln zu viel geschehen. Um zu erkennen, welche Irrationalität \gesperrt{notwendig} zu adjungieren ist, dienen die Sätze der §§ 63, 65.

Im § 63, 3. haben wir gesehen, daß die $p$te Einheitswurzel $\varrho$ rational darstellbar ist durch die Wurzeln der Teilungsgleichung und daß durch eine Substitution der Gruppe $\mathfrak{A}$, deren Determinante mit $m$ kongruent ist, $\varrho$ in $\varrho^m$ übergeht; ferner haben wir im § 65 nachgewiesen, daß durch Adjunktion sämtlicher Wurzeln einer Transformationsgleichung die Gruppe der Teilungsgleichung auf die Gruppe $\mathfrak{A}_0$ reduziert wird, die aus sämtlichen Substitutionen der Form
\weq{18}{\begin{pmatrix} a, & 0 \\ 0, & a \end{pmatrix}}
besteht, wo $a$ eine beliebige, durch $p$ nicht teilbare Zahl ist. Die Determinanten der Substitutionen von $\mathfrak{A}_0$ sind also Quadrate und sind daher nach dem Modul $p$ kongruent mit je einem der $\frac{1}{2}(p-1)$ \gesperrt{quadratischen Reste} von $p$.

Die Summe
\weq{19}{A = \overset{a}{\Sigma}\varrho^a,}
worin für $a$ die sämtlichen \gesperrt{quadratischen Reste} von $p$ zu setzen sind, bleibt daher ungeändert durch die Substitutionen
%CONT>>
